\begin{frame}{확인 문제 답}
  \small
  \begin{tabular}{@{}lp{0.74\textwidth}@{}}
    \toprule
    사례 & 답 \\
    \midrule
    (a) & \textbf{지도학습 $\cdot$ 회귀} --- 정답(소요 시간)이 붙어 있고
          예측 대상이 \textbf{연속값} \\
    (b) & \textbf{지도학습 $\cdot$ 분류} --- ``계속/중단''이 범주 라벨,
          예측이 \textbf{이산} \\
    (c) & \textbf{강화학습} --- 정답이 미리 안 주고, 가격$\cdot$조합을
          \textit{행동}으로 고르며 \textit{나중에} 돌아오는 매출 변화가
          보상. \textbf{탐험 vs 활용}이 그대로 적용 \\
    (d) & \textbf{비지도학습 $\cdot$ 군집} --- 라벨이 없고 구조만 있음 \\
    \bottomrule
  \end{tabular}
  \vskip1em
  \kq{(c)가 이번 과목의 정수다} --- ``보상이 나중에 온다'',
  ``행동이 미래를 바꾼다'', ``탐험과 활용이 충돌한다''
  세 특징이 모두 담김.
\end{frame}
